\documentclass{ltjsarticle}
\begin{document}
\section{よく使う初等関数}
以下の3つの関数をあわせて、\textbf{三角関数}と呼びます。
\[ \sin{x}, \cos{x}, \tan{x} \]
また、以下の\textbf{指数関数}や\textbf{対数関数}もよく使われます。
\[ \exp{x}, \log{x} \]
$\exp{x}$は$e^x$のことで、$e$はネイピア数です。また、$\log{x}$は$\exp{x}$の逆関数です。一般の低$a(a>0)$を持つ指数関数$a^x$の逆関数は、以下のように表します。
\[ \log_{a}{x} \]

\section{三角関数の無限級数による定義}
大学以降では、三角関数を無限級数(無限和)により定義することがあります。例えば、$\sin{x}$は
\[ \sin{x} = \lim_{n \to \infty} \sum_{k=0}^{n} \frac{(-1)^{k}} {(2k+1)!} x^{2k+1} \]
のように定義されます。
\end{document}

